% Français 9115, batch 34 — Gallica views 661–680.
% Pass: Opus 5.5 (claude-opus-5-5), 2026-09-22 — first pass, unchecked against the views by a human.
%
% What the twenty views hold. Heat, not elasticity and not number theory:
% two pieces of her own on the propagation of heat in solids, written
% after Fourier's « théorie de la chaleur », which she cites by page
% (p. 110, 141, 348) and by article (N.o 104, 111, 119, 120), and « M.r
% Fourier » by name (v. 666, 668). Ten views carry text, ten are blank.
% The leaves are single sheets hinged on grey mounting sheets, alternately
% high and low, written on the recto (and twice on the verso); the blank
% views are the backs, where the recto shows through reversed. No
% microfilm target, no repeated exposure.
%
% (1) v. 662–670: the sphere heated arbitrarily, not uniformly in each
%   concentric layer. It continues a calculation that ends on v. 660 (not
%   in this batch; glanced at): the general equation in polar coordinates
%   r, θ, φ, labelled (A), ink 324.
%   v. 662  (ink 325) the surface equation (B) of Fourier's p. 141 in
%           spherical coordinates, reduced to dv/dr + (h/k)v = 0, « de la
%           même forme » as the uniform case of p. 110; but v now varies
%           from point to point of the outer envelope.
%   v. 663  (its verso) why the surface equation keeps its form whatever
%           the motion inside; then the equation (A) obtained « d'une
%           manière spéciale », as in art. 111, from the element between
%           two cones, two meridian planes and two spheres. Ends mid-
%           sentence and runs straight on to v. 664.
%   v. 664  (ink 326) the heat through the faces of the element; the
%           general equation dv/dt = K/(r²CD){d(r² dv/dr)/dr +
%           d(sin θ dv/dθ)/(dθ sin θ) + d²v/(sin²θ dφ²)}, identical to (A);
%           a footnote argues the spherical element is proportional to
%           sin²θ (wrong: the zone varies as sin θ; kept, with a note).
%           Several letters after K are inked out solid (\ill{}).
%   v. 666  (ink 327) with θ, φ constant (A) reduces to Fourier's (A)';
%           his integral v = (1/r){a₁e^{−Kn₁²t} sin n₁r + …} (« p. 348 »)
%           still satisfies the surface equation if multiplied by
%           coefficients in θ and φ, to be chosen so as to satisfy the
%           two angular terms of (A).
%   v. 668  (ink 328) « Un cas très remarquable »: the temperatures of a
%           layer differ from point to point but the differences stay
%           unaltered during cooling; the condition is that the angular
%           part of (A) vanish.
%   v. 670  (ink 329) the example v = (1/r){…} cot θ sin φ, verified to
%           annul the angular terms; its physical meaning (heat gained
%           along a parallel compensated by « le froid » along the
%           meridian). The foot is under a large water stain; the last
%           line stops on « sur le même » and nothing in the batch
%           continues it.
% (2) v. 672–680: a new piece under a title, « Equation du mouvement varié
%   dela chaleur dans un cylindre solide », in a steadier hand, articles
%   numbered 1–4: the infinite cylinder whose sections are heated alike;
%   (A) dv/dt = K/CD (d²v/dx² + d²v/dy²), surface (B) (v. 672, ink 330);
%   polar coordinates and (C) dv/dt = K/CD {(1/r)dv/dr + d²v/dr² +
%   (1/r²)d²v/dφ²} (v. 674, ink 331); the surface equation reduced to (D)
%   dv/dr + (h/k)v = 0 and art. 3, the same « d'une manière spéciale »,
%   citing N.o 119 « (T. dela chaleur) » (v. 676, ink 332); N.o 120 and
%   art. 4, v = e^{−mt}u, u = sin βφ.r^β.s (v. 678, ink 333); the
%   equation for s solved by the series s = 1 − gr²/2(n+1) + …, n = 2β+1,
%   and the particular value e^{−gkt}(A sin βφ + B cos βφ)r^β s; the
%   page stops at « qui doit être satisfaite lorsque » (v. 680, ink 334).
%   v. 679, the verso of 333, holds only nine lines, struck with two long
%   crossing strokes: the end of an enumeration of roots g, values u and
%   arbitrary coefficients B⁰; A, B; A′, B′ … « sont des coefficiens
%   arbitraires ». It does not follow the recto; it reads like the end of
%   a later stage of the same integration, written on the back and
%   cancelled.
%
% Views not transcribed. v. 661: blank back of the leaf of v. 660 (not in
% this batch). v. 665, 667, 669, 671, 673, 675, 677: blank backs of the
% leaves 326–332, the recto showing through reversed (checked on
% thumbnails of all of them; v. 671 also shows the back of the water
% stain of v. 670 as a large dark blot).
%
% Seams. v. 660/661: view 660 (ink 324, glanced at) ends on the general
% equation (A) in r, θ, φ; view 662 takes it up as « l'équation générale
% (A) » (v. 663), so the heat-in-the-sphere piece begins before this batch.
% Batch 33, written in parallel and read after these views were
% transcribed, agrees: it opens the sphere memoir at v. 652 (« la chaleur
% dans une sphère solide échauffée d'une manière arbitraire », v. 646) and
% ends v. 660 on the same equation (A), with nothing cut at 660/661.
% v. 680/681: view 680 stops mid-sentence (« lorsque »); view 681 is the
% blank back of that leaf, so the sentence continues at v. 682 or later.
%
% The numbers on the leaves, and why no \folio{} is written. (1) The ink
% series at the top right, in the larger rounder hand of the whole
% volume: 325 (v. 662), 326 (664), 327 (666), 328 (668), 329 (670), 330
% (672), 331 (674), 332 (676), 333 (678), 334 (680); 324 on v. 660 just
% before. One per leaf, on the written recto only, never on a back (663
% and 679 are written backs and carry none). Penned with thick and thin
% strokes, as in every earlier batch: no \folio{}. (2) A second, smaller
% series of single figures beside or above the ink number, each struck
% through, in the text's ink: « 4 » preceded by an oblique stroke (662),
% « 5 » alone at the head of the verso (663, doubtful), 6 (664), 7 (666),
% none (668), 9 (670); then none (672), 2 (674, pale, crossed in black),
% none (676), 4 (678, doubtful), 5 (680, doubtful — could be 8). Read as
% it stands it looks like her own page numbers, struck when the ink series
% was added: 4–9 for the sphere piece and 1–5 for the cylinder piece, the
% missing figures (8 at 668, 1 at 672, 3 at 676) simply not written or not
% visible. This is an inference, recorded as such; batch 33 reads « 1 »
% with a small cross above the ink number at v. 652, where the sphere
% memoir opens, and a struck figure read 2 at v. 654, which fits a run
% 1, 2, … 9 across both batches. (3) Her article numbers
% 1, 2, 3, « N.o 4 » in the cylinder piece; references to Fourier's pages
% 110, 141, 348 and articles 104, 111, 119, 120 (« p. 348 » is not a folio).
% Every number above was read on its view; none is inferred.
%
% Manuscript C. Nothing here is number theory, and the ink series runs
% 325–334: Laubenbacher and Pengelley's ff. 348r–349r, if this series is
% the foliation they cite, should fall about fourteen leaves on, i.e.
% around views 706–710 at this rate of one numbered leaf per two views.
% That is an extrapolation, not a reading.
%
% Notation. Hers, kept. « sin », « cos », « cot » (with or without a stop)
% set \text{sin} etc.; K, C.D, h, k as in Fourier; she writes both K and k,
% and in the cylinder piece redefines k as K/C.D and h as h/k (v. 678).
% The capital of equation (D) is a dome on a bar, like the (D) of batch
% 28's heat draft, and could be taken for Ω. Her « &c » is set \&c.
% Letters she underlines in the prose (v, r, z, t, e, q, u, s) are set
% \underline inside the maths. Her « sin²θ » for the spherical element is
% kept (v. 664). Spelling hers, not flagged: « tems », « connoître »,
% « varieroit », « independans », « coëfficiens », « differens »,
% « suivans », « correspondans », « menerons », « egalement », « assujeti »,
% « parallele », « parconséquent », « cequi », « cequ », « laquestion »,
% « laquantité », « dela », « aulieu », « atravers », « c'est adire »,
% « a la fois », « Demême », « satisferon », « constants » (for
% constantes), « dont s'agit ».
%
% What could not be read, and what is offered with doubt. \ill{} stands
% eight times: one word of a fine interlinear addition (v. 663, « qui
% s'\ill{} sur ce parallele »); five letters inked out solid in the
% heat-flux expressions of v. 664 (after K twice, inside d( … dv/dφ), and
% before sin θ twice); two solid blots in « cot \ill{} θ sin \ill{} φ »
% (v. 670), probably the m of a struck « sin mφ ». \uncertain{} stands
% for « page 2 » (v. 662), « r² » in the footnote's last ratio (v. 664),
% « est » and the struck « couche » (v. 668), the 2 of the struck
% « 2t²−1 = m² » and the words « est v′ », « v″ », « troisième » under
% the water stain (v. 670), « les » and « mettant » in a struck line and
% the v of one dv/dφ (v. 674), and the small struck figures of v. 663,
% 678, 680. Nothing was guessed to fill a gap.
%
% Nothing here was checked against any published edition of Sophie
% Germain's papers, or against Fourier's Théorie analytique de la chaleur.
\documentclass[11pt,a4paper]{article}
% The preamble shared by every transcription of the mathematicians' manuscripts in Gallica.
%
% It rests on one idea: **uncertainty compounds, it does not resolve.** A
% transcription that smooths over a doubtful word has destroyed the only thing
% it was made for — a reader must be able to tell, at every point, what was on
% the page and what was supplied. The apparatus macros below are therefore the
% critical apparatus, not decoration.
%
% The same commands are recognised by `scripts/render.mjs`, which produces the
% site's reading view, and by `scripts/tei.mjs`. A macro added here must be
% added there too, or rendering fails loudly — which is the wanted behaviour.
%
% Comments here are English, like the rest of the repository. The documents
% this preamble sets are French, and that is deliberate: see the skills.

\ProvidesPackage{germain}[2026/09/15 Transcriptions of mathematicians' manuscripts in Gallica]

\usepackage{amsmath,amssymb,amsthm}
\usepackage{mathrsfs}
% Kept from the parent preamble so the renderer's diagram support stays
% available; these manuscripts draw no commutative diagrams, and nothing here uses it.
\usepackage{tikz-cd}
\usepackage[english,french]{babel}
\usepackage{fontspec}

% --- Greek ------------------------------------------------------------------
% The text face stays fontspec's default, Latin Modern, which has no Greek:
% XeTeX drops every glyph it lacks with a line in the log and nothing on the
% page, and the Greek words the manuscripts quote vanished from the PDFs. So
% the Greek and Greek Extended blocks get a XeTeX character class of their own,
% and entering or leaving a run of it switches the family to CMU Serif — the
% same Computer Modern design, with polytonic Greek — and back. Only the
% family changes: a Greek word in \textit or \textbf keeps its shape and series.
% The transitions to and from every other class carry the tokens that class
% has towards an ordinary letter, so French spacing before « ; » or inside
% « » still holds after a Greek word. No grouping, so no brace or \endgroup
% can come out unbalanced; maths is left alone, since XeTeX inserts nothing
% there. Kept identical in `exercices.sty`.
\makeatletter
\newfontfamily\germainGreekfont{cmun}[
  Extension=.otf, NFSSFamily=germaingreek,
  UprightFont=*rm, ItalicFont=*ti, SlantedFont=*sl,
  BoldFont=*bx, BoldItalicFont=*bi, BoldSlantedFont=*bl]
\newXeTeXintercharclass\germain@greek
\def\germain@greekrange#1#2{%
  \@tempcnta=#1\relax
  \loop
    \XeTeXcharclass\@tempcnta=\germain@greek
    \ifnum\@tempcnta<#2\advance\@tempcnta\@ne\repeat}
\germain@greekrange{"0370}{"03FF}
\germain@greekrange{"1F00}{"1FFF}
\def\germain@greekprev{\rmdefault}
\def\germain@greekin{%
  \edef\germain@greekprev{\f@family}\fontfamily{germaingreek}\selectfont}
\def\germain@greekout{\fontfamily{\germain@greekprev}\selectfont}
\def\germain@greekjoin{%
  \XeTeXinterchartokenstate=\@ne
  \@tempcnta=\z@
  \loop
    \ifnum\@tempcnta=\germain@greek\else
      \XeTeXinterchartoks\germain@greek\@tempcnta=\expandafter{%
        \expandafter\germain@greekout\the\XeTeXinterchartoks\z@\@tempcnta}%
      \XeTeXinterchartoks\@tempcnta\germain@greek=\expandafter{%
        \the\XeTeXinterchartoks\@tempcnta\z@\germain@greekin}%
    \fi
    \ifnum\@tempcnta<4095\advance\@tempcnta\@ne\repeat}
% babel-french sets its own classes, and saves and restores the interchar
% state, each time a language is selected: join again after it does.
\addto\extrasfrench{\germain@greekjoin}
\addto\extrasenglish{\germain@greekjoin}
\AtBeginDocument{\germain@greekjoin}
\makeatother

\usepackage{geometry}
\geometry{a4paper,margin=25mm,marginparwidth=28mm}
\usepackage{marginnote}
\usepackage{xcolor}
\usepackage[normalem]{ulem}
\setlength{\emergencystretch}{3em}
\usepackage{hyperref}
\hypersetup{colorlinks=true,linkcolor=black,urlcolor=black,pdfborder={0 0 0}}

% --- Batch metadata ---------------------------------------------------------
% Taken as they stand by the rendering script, which shows them at the head of
% the reading view. They fix what the file claims to cover. The macro names are
% the parent project's, so that the scripts stay shared; \folder here means the
% volume's slug — `fr-9115` — and \pages counts Gallica views.

\newcommand{\folder}[1]{\gdef\grothFolder{#1}}
\newcommand{\batch}[1]{\gdef\grothBatch{#1}}
\newcommand{\pages}[2]{\gdef\grothFirst{#1}\gdef\grothLast{#2}}
\newcommand{\foldertitle}[1]{\gdef\grothFoldertitle{#1}}

% The holder's own shelfmark, printed as they write it: « Français 9115 ».
\newcommand{\shelfmark}[1]{\gdef\grothShelfmark{#1}}

% The dating, copied verbatim from the holder's catalogue — never inferred
% here. « XIXe siècle » is the BnF's; a pass that dates a leaf from its content
% says so in a \note{}, as its own inference.
\newcommand{\dating}[1]{\gdef\grothDating{#1}}

% The status notice, printed in the title block and repeated as a diagonal
% watermark on every page of the PDF. The manuscripts are in the public
% domain; what this line declares is not a right but a quality — an
% unverified first pass by a machine. The skills write the line; a file
% without it gets no watermark, so leaving it out is visible in the output.
\newcommand{\watermark}[1]{\gdef\grothWatermark{#1}}

\gdef\grothFolder{?}\gdef\grothBatch{}\gdef\grothFirst{?}\gdef\grothLast{?}%
\gdef\grothFoldertitle{}\gdef\grothShelfmark{}\gdef\grothDating{}\gdef\grothWatermark{}

% --- The résumé -------------------------------------------------------------
\newenvironment{resume}
  {\small\setlength{\parskip}{0.45em}}
  {\par\medskip\hrule\medskip}

% --- Keywords ---------------------------------------------------------------
% In English, closing the modernised reading's résumé: the modern vocabulary
% under which to search for what the leaves construct. The manifest extracts
% them to tag the volume — this line is the single source of the tags.
\newcommand{\keywords}[1]{\par\smallskip\noindent
  {\footnotesize\textsc{Keywords} --- \textit{#1}}\par}

% --- The critical apparatus -------------------------------------------------

% The Gallica view number, in the margin. This is the anchor that lets a
% disagreement over a formula be settled by looking at *one* image rather than
% a volume of 750. The site uses it to turn the facsimile as the transcription
% is scrolled. A view is one scanned image: usually one side of a leaf.
\newcommand{\page}[1]{%
  \marginnote{\footnotesize\color{gray!70}\textsf{v.\,#1}}%
  \label{page:#1}\ignorespaces}

% The BnF's pencilled foliation, where the leaf carries it and a pass can read
% it: « 198r ». This is what the literature cites — « ff. 198r–208v » — and
% the one line that turns a citation into a view. Recorded, never inferred:
% a folio a pass cannot read is not written.
\newcommand{\folio}[1]{%
  \marginnote{\footnotesize\color{gray!50}\textsf{f.\,#1}}\ignorespaces}

% The modernised reading's coarser counterpart to \page{}: which views a
% section is drawn from.
\newcommand{\pagerange}[2]{%
  \marginnote{\footnotesize\color{gray!70}\textsf{v.\,#1--#2}}%
  \ignorespaces}

% Illegible: never guessed.
\newcommand{\ill}{\textcolor{red!60!black}{[\,\ldots\,]}}

% A doubtful reading: the word is given, and flagged as uncertain.
\newcommand{\uncertain}[1]{\uline{#1}}

% An editorial addition — a missing letter, an implied word.
\newcommand{\add}[1]{\textcolor{blue!50!black}{[#1]}}

% Struck out by the writer. What was crossed out is part of the document.
\newcommand{\struck}[1]{\sout{#1}}

% The transcriber's note, on the state of the page or the mathematics.
\newcommand{\note}[1]{\par\smallskip{\small\color{gray!60!black}\textit{[#1]}}\par\smallskip}

% A marginal note of hers, distinct from the body of the text.
\newcommand{\marginal}[1]{\par\smallskip\noindent\hspace{1em}%
  {\small\color{gray!60!black}#1}\par\smallskip}

% --- Title ------------------------------------------------------------------
% The title block is French, like the documents themselves.

\AddToHook{shipout/background}{%
  \ifx\grothWatermark\empty\else
    \begin{tikzpicture}[remember picture, overlay]
      \node[rotate=52, scale=3.4, align=center, text=gray!18,
            font=\sffamily\bfseries]
        at (current page.center) {\grothWatermark};
    \end{tikzpicture}%
  \fi}

\AtBeginDocument{%
  \begin{center}
    {\large\bfseries Manuscrits de mathématiciens (Gallica) ---
      \ifx\grothShelfmark\empty\grothFolder\else\grothShelfmark\fi}\\[2pt]
    {\itshape\grothFoldertitle}\\[4pt]
    {\small vues \grothFirst--\grothLast
      \ifx\grothBatch\empty\else\ (lot \grothBatch)\fi}\\[2pt]
    \ifx\grothDating\empty\else
      {\small\color{gray!60!black}Datation du catalogue : \grothDating}\\[2pt]
    \fi
    \ifx\grothWatermark\empty\else
      {\small\bfseries\color{red!45!gray}\grothWatermark}\\[2pt]
    \fi
    {\footnotesize\color{gray!60!black}Transcription établie d'après les images IIIF de Gallica.
     Source gallica.bnf.fr / Bibliothèque nationale de France.\\
     Les lectures incertaines sont soulignées, les illisibles notées [\,\ldots\,],
     les ajouts entre crochets ; \textsf{v.} numérote les vues de Gallica,
     \textsf{f.} la foliotation au crayon de la BnF.}
  \end{center}
  \vspace{1em}}

\endinput


\folder{fr-9115}
\shelfmark{Français 9115}
\batch{34}
\pages{661}{680}
\dating{1801-1900}
\watermark{Lecture automatique\\non vérifiée}
\foldertitle{Papiers de Sophie GERMAIN. — Recueil de dissertations et problèmes mathématiques et physiques.}

\begin{document}

\note{Numérotation des feuillets. En haut à droite, à l'encre, d'une écriture plus ronde, la série qui court dans tout le volume : 325 (vue 662), 326 (664), 327 (666), 328 (668), 329 (670), 330 (672), 331 (674), 332 (676), 333 (678), 334 (680), sur le recto seulement. Cette série occupe la place de la foliotation de la Bibliothèque, mais elle est tracée à la plume et non au crayon ; aucun « f. » n'est donc porté. À côté ou au-dessus, sur la plupart de ces feuillets, un petit chiffre barré, de l'encre du texte, relevé dans les notes de chaque vue ; ce pourrait être une pagination antérieure de sa main. Tous ces nombres ont été lus sur la vue ; aucun n'est déduit. Ne sont pas transcrites : la vue 661, dos blanc du feuillet de la vue 660 ; les vues 665, 667, 669, 671, 673, 675 et 677, dos blancs des feuillets 326 à 332, où le recto se lit par transparence.}

\page{662}

\note{En haut à droite, à l'encre, « 325 », le dernier chiffre à longue queue descendante ; au-dessus, un petit chiffre barré d'une croix, peut-être un 4, précédé d'un trait oblique. Le feuillet est monté bas sur sa feuille de garde ; il est écrit des deux côtés (vue 663).}

Pour savoir cequ devient, dans le cas présent l'équation
dela surface $m\frac{dv}{dx} + n\frac{dv}{dy} + p\frac{dv}{dz} + \frac{h}{k}vq = 0 \ldots (\text{B})$
donnée page 141 dela théorie dela chaleur il faut
\struck{déterminer} \add{connoître} les valeurs de $m$, $n$, $p$, et $q$. On voit
qu'a raison de l'équation de la sphère l'équation
$m\,dx + n\,dy + p\,dz = 0$ devient ici $x\,dx + y\,dy + z\,dz = 0$
par conséquent on a $m = x$, $n = y$, $p = z$ \quad $q = \sqrt{m^{2}+n^{2}+p^{2}} = r$
\[ = r\,\text{sin}\,\theta\,\text{cos}\,\varphi,\quad = r\,\text{sin}\,\theta\,\text{sin}\,\varphi,\quad = r\,\text{cos}\,\theta \]
\note{« cequ » : ainsi, sans « e » final. « connoître » est écrit au-dessus de « déterminer », souligné en guise de rature. La troisième ligne de valeurs est écrite sous $m = x$, $n = y$, $p = z$, terme à terme.}

ces valeurs mises dans l'équation (B) en même tems que celle
de $\frac{dv}{dx}$, $\frac{dv}{dy}$, $\frac{dv}{dz}$; trouvées page \uncertain{2} donnent
\note{Le chiffre après « page » a la forme de son $z$ ; lu 2 avec doute.}
\[ r\,\text{sin.}\,\theta\,\text{cos}\,\varphi\left\{\frac{dv}{dr}.\text{sin}\,\theta\,\text{cos}\,\varphi + \frac{dv}{d\theta}\,\frac{\text{cos.}\,\theta\,\text{cos}\,\varphi}{r} - \frac{dv}{d\varphi}.\frac{\text{sin.}\,\varphi}{r\,\text{sin}\,\theta}\right\} \]
\[ + r\,\text{sin.}\,\theta\,\text{sin.}\,\varphi\left\{\frac{dv}{dr}\,\text{sin.}\,\theta\,\text{sin}\,\varphi + \frac{dv}{d\theta}.\frac{\text{cos}\,\theta\,\text{sin.}\,\varphi}{r} + \frac{dv}{d\varphi}\,\frac{\text{cos}\,\varphi}{r\,\text{sin.}\,\theta}\right\} \]
\[ + r\,\text{cos.}\,\theta\left\{\frac{dv}{dr}\,\text{cos.}\,\theta - \frac{dv}{d\theta}.\frac{\text{sin}\,\theta}{r}\right\} + \frac{h}{k}\,v\,r = 0 \]

En effaçant cequi se détruit et en réduisant cette équation
devient \quad $\frac{dv}{dr} + \frac{h}{k}\,v = 0$

Elle est dela même forme que celle qui a été donnée p
110 et qui appartient au cas ou la chaleur est uniformement
répandue dans chacune des couches concentriques dela sphère

Il faut observer seulement qu'ici la chaleur des differens
points de l'enveloppe qui est exposée au contact de l'air
n'étant pas la même la valeur de $\underline{v}$ differe d'un point
a l'autre. Les coëfficiens qui doivent entrer dans l'intégrale
de l'équation dela sphère échauffée d'une manière arbitraire
se réduiroient a des constantes si la chaleur étoit uniformement
répandue dans chacune des couches, ainsi la valeur de $\underline{v}$ qui
doit être variable ici deviendroit alors constante. ainsi
\note{Le dernier mot, « ainsi », se poursuit au verso (vue 663).}

\page{663}

\note{Le verso du feuillet 325 (vue 662), écrit à la suite. En haut à droite, à la fin de la première ligne, un petit chiffre isolé, lu \uncertain{5}, qui ferait suite au chiffre barré du recto ; pas de nombre de la série à l'encre. Au-dessus du feuillet, sur la feuille de montage, l'écriture d'un autre feuillet paraît à l'envers par transparence.}

l'analyse montre aussi bien que le simple raisonnement que
l'équation dela surface conserve la même forme soit que
la chaleur soit uniformement répandue dans tous les
points également eloignés dela sphère soit que ces
\struck{differens} points \struck{ayent des tem} soient élevés a des
temperatures differentes; il est evident en effet qu'a
la surface le mouvement ne peut s'exécuter que dans
une seule direction et que par conséquent soit qu'a l'égard
\marginal{des points} \add{intérieurs} \struck{qu} on considère le mouvement dans trois, deux,
ou une seule dimension. l'équation qui représente
le mouvement \add{a la surface} est nécessairement dela même forme.
\note{« aussi » est récrit en traits appuyés, peut-être sur un autre mot. « des points » est écrit dans la marge de gauche, devant la ligne, et « intérieurs » au-dessus de « qu'on », dont « qu » est barré : la phrase se lit « soit qu'a l'égard des points intérieurs on considère ». « dimension » : ainsi, au singulier.}

On obtiendroit également l'équation générale (A)
en traitant laquestion d'une manière spéciale
lorsqu'il s'est agi art 111 et suivans de trouver l'équation
du mouvement de la chaleur dans une sphère solide
dont tous les points également eloignés du centre ont une
temperature uniforme on a pû considerer a la fois
toute l'étendue $4\pi r^{2}$ dela surface. Dans le cas présent
ou la chaleur se meut a la fois dans plusieurs directions
il faut considerer l'élément de la couche comprise
entre les deux surfaces concentriques dont les rayons
sont $r$ et $r + dr$. Pour determiner la forme de cette
élément nous concevrons une \struck{surface} conique dont le
sommet soit au centre dela sphère et qui ait
pour base un des parallele. \add{qui s'\ill{} sur ce parallele} \add{et soit terminé par la calotte spherique} par exemple celui dont
la latitude est $\theta$. nous concevrons en même tems
un second cône \struck{dont le} concentrique au premier et qui
ait pour base le parallele dont la latitude est $\theta + d\theta$
\note{Deux additions superposées entre les lignes, au-dessus de « par exemple » : la plus basse, d'une plume très fine, « qui s'… sur ce parallele », dont le verbe n'a pu être lu ; la plus haute, « et soit terminé par la calotte spherique », qui semble la remplacer. Rien n'indique sur le feuillet l'ordre où elles s'insèrent. « parallele » après « un des » : ainsi, au singulier. La page s'arrête au bas du feuillet ; la suite n'est pas sur la vue suivante, qui est un autre feuillet (voir vue 664).}

\page{664}

\note{En haut à droite, à l'encre, « 326 » ; au-dessus, un petit « 6 » barré d'un trait, comme le chiffre barré de la vue 662. Le feuillet est monté haut sur sa feuille de garde. Le texte continue la vue 663 sans rupture : « … dont la latitude est $\theta + d\theta$ » / « enfin nous menerons les deux plans … ».}

enfin nous menerons les deux plans correspondans aux longitudes
$\varphi$ et $\varphi + d\varphi$. nous considererons donc le mouvement dela chaleur
dans l'élément sphérique compris entre les deux surfaces coniques dont
les bases sont les cercles dont les rayons sont $r\,\text{sin}\,\theta$ et $r(\text{sin}\,\theta + d\theta)$
les deux plans dont les longitudes sont $\varphi$ et $\varphi + d\varphi$ et les deux
couches concentriques dont les rayons sont $r$ et $r + dr$. le point
$m$ appartiendra aux valeurs $r\,\text{sin}\,\theta$, $r$ et $\varphi$ et le point $m'$ aux
valeurs \struck{$r\,\text{sin}\,\theta + d\theta$} $(r + dr)(\text{sin}\,\theta + d\theta)$, $r + dr$ et $\varphi + d\varphi$
\note{Dans « $r(\text{sin}\,\theta + d\theta)$ », le second $\theta$ est récrit sur une autre lettre et taché ; la parenthèse est fermée au bord du feuillet. Ainsi écrit, sans parenthèse autour de $\theta + d\theta$.}

La quantité de chaleur qui pénètre pendant l'instant $dt$ dans la molécule atravers
l'élément $r^{2}\,\text{sin}^{2}\theta\,d\theta\,d\varphi$ * de la surface sphérique perpendiculaire a $r$
est $-K\,r^{2}\,\text{sin}^{2}\theta\,d\theta\,d\varphi.\frac{dv}{dr}\,dt$, \add{pour avoir} celle qui sort en même tems dela molecule
par la face opposée c'est adire par l'élément dela surface sphérique dont le rayon
est $r + dr$, il faut ajouter a $-K\,r^{2}\,\text{sin}^{2}\theta\,d\theta\,d\varphi.\frac{dv}{dr}.dt$ laquantité
$-K\,\text{sin}^{2}\theta\,d\theta\,d\varphi\,d\left(r^{2}.\frac{dv}{dr}\right)dt$.
\note{L'astérisque après $d\varphi$ renvoie à la note du bas de la page. L'exposant de « sin » est le même petit 2 que celui de $r^{2}$, à chaque occurrence de ce feuillet, et la note du bas le confirme (« proportionnelle a sin$^{2}\theta$ ») ; l'élément de la sphère est en fait $r^{2}\,\text{sin}\,\theta\,d\theta\,d\varphi$. La lettre $v$ du premier $\frac{dv}{dr}$ est récrite.}

Laquantité de chaleur qui pendant le même instant $dt$ entre dans \struck{la mol} l'élément
sphérique atravers la surface conique comprise entre les deux plans
dont les longitudes sont $\varphi$ et $\varphi + d\varphi$ est $-K\,\ill{}\,d\theta\,dr\,\frac{dv}{d\varphi}$
La differentielle $-K\,\ill{}\,d\theta\,dr\,d\left(\ill{}\,\frac{dv}{d\varphi}\right)$ prise par rapport a $\varphi$
seulement est laquantité de chaleur que \struck{l'élément} l'élément acquiert
\note{Après chaque $K$ de ces deux lignes, et devant $\frac{dv}{d\varphi}$ dans la parenthèse, des lettres noircies à l'encre jusqu'à former des taches pleines : effacements de sa main, qu'on ne peut lire. La lettre lue $d\theta$ est récrite sur un $d\varphi$. « rapport » est écrit « rapports ».}

*( on trouve $\text{sin}\,\theta\,d\varphi$ parceque $r\,\text{sin}\,\theta$ est le rayon du parallele)
\note{Ligne serrée entre les autres, commençant dans la marge par un astérisque et une parenthèse ; elle glose, semble-t-il, le facteur $\text{sin}\,\theta\,d\varphi$ des expressions précédentes.}

Enfin laquantité de chaleur qui passe par la section cylindrique
comprise entre les deux plans dont les latitudes sont $\theta$ et $\theta + d\theta$
section dont la surface est \struck{$dr\,d\varphi$ est $-K$} $\ill{}\,\text{sin}\,\theta\,d\varphi\,dr$ *( )
est $-K\,\ill{}\,\text{sin}\,\theta\,d\varphi\,dr\,\frac{dv}{d\theta}$ et en prenant la differentielle de cette
quantité par rapport a $\theta$ seulement on a la chaleur acquise
par l'élément sphérique dans la direction de la latitude $\theta$.
\note{Devant $\text{sin}\,\theta$, aux deux lignes, la même tache pleine qu'après $K$ plus haut. Après $dr$, un astérisque suivi de deux parenthèses vides : un renvoi resté sans texte.}

On trouve donc que l'équation generale du mouvement de
la chaleur est.
\[ \frac{dv}{dt} = \frac{K}{r^{2}.C.D}\left\{\frac{d\left\{r^{2}\frac{dv}{dr}\right\}}{dr} + \frac{d\left(\text{sin}\,\theta.\frac{dv}{d\theta}\right)}{d\theta\,\text{sin}\,\theta} + \frac{d^{2}v}{\text{sin}^{2}\theta\,d\varphi^{2}}\right\} \]
\note{L'exposant de $r^{2}$ dans la première accolade est couvert d'une tache ; le dénominateur $dr$ du premier terme est noirci et récrit. Au second terme, « sin $\theta$ » est écrit sous le trait, au-dessous de $d\theta$. Au troisième, le dénominateur est récrit : « sin » suivi d'un groupe surchargé, puis « ($\varphi$ 2) » ; lu $\text{sin}^{2}\theta\,d\varphi^{2}$ avec doute.}

Les differens termes de cette équation étant développés
\struck{donnent au} le second membre est identiquement le même que celui de l'équation (A)

\note{Un trait horizontal traverse ici le feuillet et sépare le texte de la note du bas.}

* Il est facile de se convaincre que l'élément de la surface sphérique est en effet
proportionnelle a $\text{sin}^{2}\theta$ car si on conçoit le demi fuseau terminé a l'équateur la surface
\struck{surface} sera proportionnelle a celle du segment terminé sur le plan de l'équateur par les deux rayons
qui sont les projections des deux grands arcs qui terminent le fuseau. Or si on cherche de la surface dela
partie du demi fuseau terminé au parallele de 45$^{\text{d}}$ delatitude on aura donc cette surface a la surface entière $::\ \frac{1}{2}r^{2} : r^{2}$
$::\ \uncertain{r^{2}}\,\text{sin}^{2}\theta : r^{2}$.
\note{La note est d'une écriture plus petite, au bas du feuillet. Sa conclusion est fausse : l'aire de la zone sphérique varie comme $\text{sin}\,\theta$ (ou $1 - \text{cos}\,\theta$ selon l'origine de la latitude), non comme $\text{sin}^{2}\theta$ ; c'est l'aire de la projection sur l'équateur qui suit $\text{sin}^{2}\theta$. La transcription garde ce qui est écrit. Le dernier rapport est écrit en petit sous le précédent.}

\page{666}

\note{En haut à droite, à l'encre, « 327 », puis, un peu plus à droite, un « 7 » barré d'un trait oblique. Le feuillet est monté bas ; le bas du feuillet est blanc. Le texte continue la vue 664 (le verso du feuillet 326, vue 665, est blanc).}

Il est évident que si on \struck{multiplie} fait $\varphi$ et $\theta$ constans c'est
adire si on admet que le mouvement dela chaleur se fait
uniquement dans le sens des $\underline{r}$ l'équation (A) se réduira
a \quad $\frac{dv}{dt} = \frac{K}{C.D}\left\{\frac{2}{r}\frac{dv}{dr} + \frac{d^{2}v}{dr^{2}}\right\} \ldots\ldots (\text{A})'$

qui est en effet celle qui a été donnée par M$^{\text{r}}$ Fourier.

Nous avons vu que soit qu'on se borne a l'examen de ce cas
particulier soit qu'on considère le cas général \struck{l'équation
dela surface} on arrive a la même équation relativement
a la surface. \add{mais que dans le cas général les valeurs de $\underline{v}$ different d'un point a un autre} Il résulte dela que l'intégrale
\[ v = \frac{1}{r}\left\{a_{1}e^{-Kn_{1}^{2}t}\,\text{sin}\,n_{1}r + a_{2}e^{-Kn_{2}^{2}t}\,\text{sin.}\,n_{2}r + a_{3}e^{-Kn_{3}^{2}t}\,\text{sin}\,n_{3}r + \&c\right\} \]
qui a été donnée p. 348 dela théorie dela chaleur
satisfera encore a l'équation dela surface si le second
membre est multiplié par des coëfficiens quelconques
independans de $r$. Pour satisfaire a la fois a cette
equation et a l'équation (A) il ne s'agit donc
que de chercher des coëfficiens en $\theta$ et $\varphi$ \struck{qui}
qui en même tems qu'ils \struck{qu'ils} donneront pour
$\underline{v}$ une valeur differente pour les differens points
dela surface pour les quels les \struck{les} valeurs de $\theta$ et
$\varphi$ ne sont pas les mêmes; satisferon a l'existence
des deux termes $\frac{K}{r^{2}C.D}\left\{\frac{d\left(\text{sin}\,\theta\,\frac{dv}{d\theta}\right)}{d\theta\;\text{sin.}\,\theta} + \frac{d^{2}v}{\text{sin}^{2}\theta\,d\varphi^{2}}\right\}$
de l'équation (A)
\note{« multiplie » est souligné d'un trait qui le raye. L'addition « mais que dans le cas général … d'un point a un autre » est écrite au-dessus de « résulte dela que l'intégrale », sa fin rejetée au-dessus du bord droit ; elle est placée où elle se lit, avant « Il résulte ». Dans l'intégrale, l'indice du premier $a$ est écrit devant la lettre ; « \&c » est bouclé. « p. 348 » est le renvoi d'une page imprimée de la théorie de la chaleur, non un numéro de feuillet. « satisferon » : ainsi, sans t. Au second terme de la dernière formule, « sin.$\theta$ » est écrit sous le trait, au-dessous de $d\theta$.}

\page{668}

\note{En haut à droite, à l'encre, « 328 », sans autre chiffre à côté. Le feuillet est monté haut ; le tiers inférieur est blanc. Le texte ne reprend pas une phrase de la vue 666 : il ouvre un nouvel alinéa.}

Un cas très remarquable est celui dans lequel la valeur de
$\underline{v}$ ne varieroit avec le tems qu'a raison du mouvement de
la chaleur dans le sens des $\underline{r}$ ensorte que quoique
cette valeur fut differente dans chacun des point également
distans du centre dela sphère les changemens de temperature
de chacun de ces points seroit indépendans de ces
differences qui resteroient aussi inalterables pendant la
durée du refroidissement que l'\uncertain{est} \struck{la \uncertain{couche}}
l'uniformité de temperature entre tous les points
d'une même couche sphérique lorsqu'il s'agit du
cas particulier qu'a examiné M$^{\text{r}}$ Fourier.
\note{« resteroient » : la terminaison « ent » est récrite sur un « t ». Le mot lu « l'est » porte un point qui le fait ressembler à « l'ait ». Le mot barré après lui est à demi couvert par le trait.}

Pour que la condition dont je viens de parler soit
remplie il faut que le coëfficient de $\underline{t}$ dans
l'exposant de $\underline{e}$ soit independans des coëfficiens
en $\theta$ et $\varphi$ qui \struck{sont} satisfont a l'existence
des termes $\frac{K}{r^{2}C.D}\left\{\frac{d\left(\text{sin.}\,\theta\,\frac{dv}{d\theta}\right)}{d\theta\;\text{sin.}\,\theta} + \frac{d^{2}v}{\text{sin}^{2}\theta\,d\varphi^{2}}\right\}$ dans
l'équation (A).

Il est évident que si ces coëfficiens sont tels qu'on
ait toujours \quad $\frac{d\left(\text{sin.}\,\theta.\frac{dv}{d\theta}\right)}{d\theta\;\text{sin.}\,\theta} + \frac{d^{2}v}{\text{sin}^{2}\theta\,d\varphi^{2}} = 0$

la condition dont il s'agit \add{sera remplie}, car l'équation
(A) se réduira réellement alors a la forme (A)$'$
\note{Comme aux vues 664 et 666, « sin.$\theta$ » est écrit sous le trait, au-dessous de $d\theta$.}

\page{670}

\note{En haut à droite, à l'encre, « 329 » ; à sa gauche, un peu plus haut, un « 9 » barré d'un trait. Le feuillet est monté bas ; une large tache d'eau couvre le bas de la page et assombrit les quatre dernières lignes. Le texte fait suite à la vue 668 (« a la forme (A)$'$ » / « si on prend … »).}

\struck{si on prend $\uncertain{2}t^{2} - 1 = m^{2}$ et qu'on fasse}
\[ v = \frac{1}{r}\left\{a_{1}e^{-Kn_{1}^{2}t}\,\text{sin.}\,n_{1}r + a_{2}e^{-Kn_{2}^{2}t}\,\text{sin.}\,n_{2}r + a_{3}e^{-Kn_{3}^{2}t}\,\text{sin.}\,n_{3}r\right\}\text{cot}\,\ill{}\,\theta\;\text{sin}\,\ill{}\,\varphi \]
les conditions dont s'agit seront remplies.
\note{La première ligne est barrée d'un trait. Le chiffre devant $t^{2}$ a la forme de son $z$ et du 2 de la vue 662 ; lu 2 avec doute. Devant $\theta$ et devant $\varphi$, une lettre noircie jusqu'à la tache pleine, comme à la vue 664 : peut-être le $m$ du « sin m$\varphi$ » barré à la ligne suivante, effacé ; la vérification qui suit porte sur $\text{cot}\,\theta\,\text{sin}\,\varphi$. « dont s'agit » : ainsi, sans « il ».}

\struck{En effet $\frac{d\left\{\text{sin}\,\theta\,d\left\{\frac{\text{cot}\,(\theta\,\text{sin}\,m\varphi)}{d\theta}\right\}\right\}}{\text{sin}\,\theta\,d\theta} + \frac{d^{2}\left\{\text{cot.}\,(\theta\,\text{sin}\,m\varphi)\right\}}{\text{sin}^{2}\theta\,d\varphi^{2}}$}
\note{Première rédaction de la vérification, barrée d'un long trait. Les parenthèses sont placées ainsi sur le feuillet ; le $m$ de « sin m$\varphi$ » est récrit.}

\[ \text{En effet}\quad \frac{d\left(\text{sin}\,\theta.\,d\left(\frac{\text{cot.}\,\theta\;\text{sin}\,\varphi}{d\theta}\right)\right)}{\text{sin}\,\theta\;d\theta} + \frac{d^{2}(\text{cot}\,\theta\;\text{sin}\,\varphi)}{\text{sin}^{2}\theta\;d\varphi^{2}} \]
\[ = -\,\frac{d\left(\text{sin}\,\theta.\frac{1}{\text{sin}^{2}\theta}\right)}{\text{sin}\,\theta\;d\theta}\,\text{sin}\,\varphi - \frac{\text{cot}\,\theta}{\text{sin}^{2}\theta}.\,\text{sin}\,\varphi \]
\[ = \text{sin}\,\varphi.\,\frac{\text{cos}\,\theta}{\text{sin}^{3}\theta} - \frac{\text{cot}\,\theta}{\text{sin}^{2}\theta}\,\text{sin}\,\varphi = 0. \]
\note{Au-dessus de la première ligne, « = sin m$\varphi$ », barré. Dans « cot.$\theta$ sin$\varphi$ », le $\varphi$ est récrit sur un autre signe. Le calcul est juste : $\text{cot}\,\theta\,\text{sin}\,\varphi$ annule les deux termes.}

Le cas dont il s'agit est celui ou en vertu dela
difference des temperatures \add{des divers points} d'une même enveloppe sphérique
la quantité de chaleur qui passe d'une molecule plus
échauffée que celle qui la touche sur le même parallele
dans cette dernière molécule et devroit par conséquent
élever la temperature dela même quantité est exactement
compensée par le froid qui passe dela molécule située
aussi a côté dela seconde molécule mais sur le même
méridien.

Si on nomme $\underline{v}$ la temperature dela molécule
qui donne une certaine quantité de chaleur a la molécule
située a côté dela première sur le même parallele, et dont
la temperature \uncertain{est $v'$}, il faut que \uncertain{$v''$} étant la tempera
ture d'une \uncertain{troisième} molécule située sur le même
\note{Ces quatre lignes sont sous la tache d'eau ; les mots offerts avec doute y sont à peine visibles. La lettre lue $v''$ est un $v$ bouclé, dont les accents ne se distinguent pas. La page s'arrête sur « le même », au bas du feuillet ; le verso (vue 671) est blanc et la vue 672 ouvre un autre sujet : la phrase n'a pas de suite dans ce lot.}

\page{672}

\note{En haut à droite, à l'encre, « 330 », seul. Le feuillet est monté haut. Une autre pièce commence ici, sous un titre, dans une écriture plus droite et plus posée que celle des vues 662–670, presque sans rature.}

Equation du mouvement varié dela chaleur dans un
cylindre solide

1. Supposons que toutes les tranches paralleles a la base d'un
cylindre solide d'une longueur infinie et dont le côté est
perpendiculaire a la base circulaire, ayant été échauffés
d'une manière entièrement arbitraire, mais semblable pour
chacune d'elles; le cylindre soit exposé ensuite à un
courant d'air plus froid.

Il s'agit de trouver les équations qui doivent représenter
le mouvement dela chaleur dans un tel solide.

On voit d'abord que l'équation générale
\[ \frac{dv}{dt} = \frac{K}{C.D}\left(\frac{d^{2}v}{dx^{2}} + \frac{d^{2}v}{dy^{2}} + \frac{d^{2}v}{dz^{2}}\right) \]
qui sert a déterminer le mouvement dela chaleur
dans l'interieur des corps solides se réduit ici à
\[ \frac{dv}{dt} = \frac{K}{C.D}\left(\frac{d^{2}v}{dx^{2}} + \frac{d^{2}v}{dy^{2}}\right) \ldots\ldots (\text{A}) \]
car, suivant l'hypothèse les tranches du cylindre
étant échauffés d'une manière semblable, tous les points
qui sont situés sur des lignes paralleles a \struck{l'axe du
cylindre} \add{son axe} qui est en même tems celui des $\underline{z}$, ont la
la même temperature, parconséquent le mouvement
dela chaleur est nulle dans cette direction, et on
a \quad $\frac{d^{2}v}{dz^{2}} = 0$

Par la même raison l'équation générale relative a la
surface se réduit à \qquad $m\,\frac{dv}{dx} + n\,\frac{dv}{dy} + \frac{h}{k}\,vq = 0 \ldots (\text{B})$
\note{« son axe » est écrit au-dessus de « l'axe du cylindre », barré. « la la même » : l'article est répété d'une ligne à l'autre. Le bas du feuillet est blanc sous l'équation (B).}

\page{674}

\note{En haut à droite, à l'encre, « 331 » ; à sa droite, un petit « 2 » d'une encre pâle, barré d'un trait noir oblique. Le feuillet est monté bas. Le texte suit la vue 672 (le verso, vue 673, est blanc) : article 2 après l'article 1.}

2. Le centre d'une des tranches paralleles a la base du
cylindre étant pris pour origine; soit $\underline{r}$ le rayon du
cercle qui mesure une des couches de cette tranche et $\varphi$
l'arc compris entre un des points dela circonference
et celui auquel appartient l'ordonnée $y = 0$, on aura
quel que soit l'arc $\varphi$.
\[ x = r\,\text{cos.}\,\varphi,\quad y = r\,\text{sin.}\,\varphi; \]
\[ dx = dr\,\text{cos}\,\varphi - r\,d\varphi\,\text{sin.}\,\varphi,\qquad dy = dr\,\text{sin.}\,\varphi + r\,d\varphi\,\text{cos}\,\varphi \]
\struck{$\frac{dx}{dr} = \text{cos}\,\varphi,\quad \frac{dy}{dr} = \text{sin.}\,\varphi$}
\[ dr = dx\,\text{cos.}\,\varphi + dy\,\text{sin.}\,\varphi,\qquad d\varphi = \frac{dy\,\text{cos.}\,\varphi - dx\,\text{sin.}\,\varphi}{r} \]
\[ \frac{dr}{dx} = \text{cos.}\,\varphi,\quad \frac{dr}{dy} = \text{sin.}\,\varphi;\quad \frac{d\varphi}{dx} = -\,\frac{\text{sin}\,\varphi}{r},\quad \frac{d\varphi}{dy} = \frac{\text{cos.}\,\varphi}{r} \]

\struck{En substituant \uncertain{les} $v$ \uncertain{mettant} aulieu}.
ces valeurs de \struck{$\frac{dx}{dy}$} $\frac{dr}{dx}$, $\frac{dr}{dy}$, $\frac{d\varphi}{dx}$ et $\frac{d\varphi}{dy}$ étant substituées
dans celles de $\frac{dv}{dx}$, $\frac{dv}{dy}$, $\frac{d^{2}v}{dx^{2}}$ et $\frac{d^{2}v}{dy^{2}}$ on trouve
\[ \frac{dv}{dx} = \text{cos.}\,\varphi.\,\frac{dv}{dr} - \frac{\text{sin.}\,\varphi}{r}.\,\frac{dv}{d\varphi} \]
\[ \frac{dv}{dy} = \text{sin.}\,\varphi.\,\frac{dv}{dr} + \frac{\text{cos.}\,\varphi}{r}.\,\frac{dv}{d\varphi} \]
\[ \frac{d^{2}v}{dx^{2}} = \frac{\text{sin}^{2}\varphi}{r}.\,\frac{dv}{dr} + 2\,\frac{\text{sin.}\,\varphi\,\text{cos.}\,\varphi}{r^{2}}.\,\frac{d\uncertain{v}}{d\varphi} + \text{cos}^{2}\varphi.\,\frac{d^{2}v}{dr^{2}} - 2\,\frac{\text{sin.}\,\varphi\,\text{cos.}\,\varphi}{r}.\,\frac{d^{2}v}{dr\,d\varphi} \]
\[ + \frac{\text{sin}^{2}\varphi}{r^{2}}.\,\frac{d^{2}v}{d\varphi^{2}} \]
\[ \frac{d^{2}v}{dy^{2}} = \frac{\text{cos}^{2}\varphi}{r}.\,\frac{dv}{dr} - 2\,\frac{\text{sin.}\,\varphi\,\text{cos.}\,\varphi}{r^{2}}.\,\frac{dv}{d\varphi} + \text{sin}^{2}\varphi\,\frac{d^{2}v}{dr^{2}} + 2\,\frac{\text{sin.}\,\varphi\,\text{cos}\,\varphi}{r}\,\frac{d^{2}v}{dr\,d\varphi} \]
\[ + \frac{\text{cos}^{2}\varphi}{r^{2}}.\,\frac{d^{2}v}{d\varphi^{2}} \]
\[ \frac{d^{2}v}{dx^{2}} + \frac{d^{2}v}{dy^{2}} = \frac{1}{r}\,\frac{dv}{dr} + \frac{d^{2}v}{dr^{2}} + \frac{1}{r^{2}}\,\frac{d^{2}v}{d\varphi^{2}} \]
et l'équation (A) devient
\[ \frac{dv}{dt} = \frac{K}{C.D}\left\{\frac{1}{r}\,\frac{dv}{dr} + \frac{d^{2}v}{dr^{2}} + \frac{1}{r^{2}}\,\frac{d^{2}v}{d\varphi^{2}}\right\} \ldots\ldots (\text{C}) \]
\note{La ligne « En substituant … aulieu » est barrée d'un trait ; deux de ses mots sont couverts par le trait. Devant $\frac{dr}{dx}$, une première fraction, lue $\frac{dx}{dy}$, est barrée. Dans $\frac{d^{2}v}{dx^{2}}$, la lettre du numérateur de $\frac{dv}{d\varphi}$ ressemble à un $\varphi$ ; dans $\frac{d^{2}v}{dy^{2}}$, le $v$ du même terme est récrit sur un $x$, le « cos » du premier terme sur « sin », et le « sin » du troisième sur « cos ». Le premier terme de chacune des deux dérivées secondes s'écrit, sur le feuillet, avec $r$ seul au dénominateur ; la somme est juste.}

\page{676}

\note{En haut à droite, à l'encre, « 332 », un trait tracé au-dessus du nombre. Pas d'autre chiffre. Le feuillet est monté haut. Le texte suit la vue 674 (le verso, vue 675, est blanc).}

A l'égard de l'équation relative a la surface il faut remarquer
que la condition $m\,dx + n\,dy = 0$ qui donne les valeurs de
$m$ et $n$, et parconséquent aussi celle de $\underline{q}$, étant, ici, $x\,dx + y\,dy = 0$;
on peut mettre aulieu de $m$ $n$ et $q$; $x$, $y$ et $\sqrt{x^{2}+y^{2}}$
ou plutôt leurs valeurs $r\,\text{cos}\,\varphi$, $r\,\text{sin.}\,\varphi$, et $\underline{r}$. Si on \struck{met}
met aussi a la place de $\frac{dv}{dx}$ et $\frac{dv}{dy}$ les valeurs trouvées
plus haut, l'équation (B) deviendra:
\[ r\,\text{cos.}\,\varphi\left(\text{cos}\,\varphi\,\frac{dv}{dr} - \frac{\text{sin.}\,\varphi}{r}.\,\frac{dv}{d\varphi}\right) + r\,\text{sin.}\,\varphi\left(\text{sin.}\,\varphi\,\frac{dv}{dr} + \frac{\text{cos}\,\varphi}{r}\,\frac{dv}{d\varphi}\right) + r\,\frac{h}{k}\,v = 0 \]
\struck{elle} se réduira donc à
\[ \frac{dv}{dr} + \frac{h}{k}.\,v = 0 \ldots\ldots (\text{D}) \]

3. En traitant la question d'une manière spéciale
on obtiendroit également les deux équations (C) et (D).
Pour y parvenir il faut observer que dans l'hypothèse
présente le mouvement dela chaleur a lieu a la fois
dans deux directions differentes; car elle pénètre dans la
tranche infiniment petite dela couche du cylindre qui
\struck{appartient a un disque} \add{elle même est} d'une epaisseur aussi infiniment
petite.
\note{Le « (D) » des équations est une capitale en dôme sur un trait, qui ressemble à un $\Omega$ ; la suite (A), (B), (C), (D) la fait lire D. « deux » et « faut » sont traversés d'un trait horizontal qui paraît être son trait de liaison plutôt qu'une rature : le sens les demande l'un et l'autre. « elle même est » est écrit au-dessus de « a un disque », barré avec « appartient ».}

En appliquant ici l'analyse employée N.o 119 \add{(T. dela chaleur)}, on voit
que l'accroissement dela chaleur qui pendant l'instant
$dt$ est dû au mouvement qui s'exécute dans la couche
comprise entre la surface dont le rayon est $\underline{r}$ et
celle dont le rayon est $r + dr$ fournit dans l'équation
(C) les deux termes $\frac{K}{C.D}\left(\frac{d^{2}v}{dr^{2}} + \frac{1}{r}\,\frac{dv}{dr}\right)$; Demême
aussi l'analyse employée N.o 104 peut servir a prouver
\note{« (T. dela chaleur) » est écrit en petit au-dessus de « N.o 119 », entre parenthèses : le renvoi est à un article numéroté de la théorie de la chaleur. Dans « N.o 104 », le 1 est un trait court et le 4 a la forme d'un h, comme ailleurs dans sa main. La page s'arrête au bas du feuillet sur « a prouver » ; la phrase se poursuit à la vue 678.}

\page{678}

\note{En haut à droite, à l'encre, « 333 » ; à sa gauche, un petit chiffre barré d'une croix, lu \uncertain{4}. Le feuillet est monté bas ; il est écrit aussi au verso (vue 679). Le texte suit la vue 676 sans rupture : « peut servir a prouver » / « que le troisième terme … ».}

que le troisième terme $\frac{K}{C.D}.\,\frac{1}{rr}\,\frac{d^{2}v}{d\varphi^{2}}$ \struck{que le troisième terme}
du second membre dela même équation (C) est dû au
mouvement dela chaleur qui pendant l'instant $dt$
s'exécute dans la tranche infiniment petite comprise entre
une section placée a la distance $r\varphi$ et une autre section
placée a la distance $r\varphi + r\,d\varphi$.

A l'égard de l'équation (D) elle est absolument la même
que celle qui a été donnée N.o 120, il faut observer
seulement que dans le cas du cylindre uniformement
échauffé la valeur de $\underline{v}$ est la même dans tous les
points extérieurs du solide tandis qu'ici elle est differente
pour chacun d'eux.

N.o 4 Pour intégrer les équations
\[ \frac{dv}{dt} = \frac{K}{C.D}\left\{\frac{1}{r}\,\frac{dv}{dr} + \frac{d^{2}v}{dr^{2}} + \frac{1}{r^{2}}\,\frac{d^{2}v}{d\varphi^{2}}\right\},\qquad \frac{dv}{dr} + \frac{h}{k} = 0 \]
on donnera en premier lieu a $\underline{v}$ une valeur particulière
exprimée par l'équation $v = e^{-mt}u$; $m$ est un nombre
quelconque et $\underline{u}$ une fonction de $\underline{r}$ et de $\varphi$. \struck{e}
On designe par $k$ le coëfficient \add{$\frac{K}{C.D}$ qui entre dans la première equation et par $h$ le coëfficient} $\frac{h}{k}$ qui entre dans la seconde.
En substituant la valeur attribuée a $\underline{v}$ on trouve la condition
suivante:
\[ \frac{m}{k}\,u + \frac{d^{2}u}{dr^{2}} + \frac{1}{r}\,\frac{du}{dr} + \frac{1}{r^{2}}.\,\frac{d^{2}u}{d\varphi^{2}} = 0 \]
On donnera ensuite a $\underline{u}$ la valeur particulière exprimée
par l'équation \quad $u = \text{sin}\,\beta\varphi.\,r^{\beta}.\,s$ ; $\beta$ est un nombre
entier et $s$ une fonction de $\underline{r}$ seule.
\note{« N.o 4 » est écrit dans la marge de gauche, devant la ligne : c'est le numéro de l'article, après les articles 1 à 3 des vues 672–676. Dans la seconde équation, $v$ manque après $\frac{h}{k}$ ; ainsi sur le feuillet. L'addition « $\frac{K}{C.D}$ qui entre dans la première equation et par $h$ le coëfficient » est écrite entre les lignes, au-dessus de « le coëfficient $\frac{h}{k}$ qui entre dans la seconde », sa fin rejetée au-dessus du bord droit ; elle se place devant $\frac{h}{k}$. « par » devant $k$ est traversé d'un trait, peut-être de liaison. Le bas du feuillet est blanc.}

\page{679}

\note{Le verso du feuillet 333 (vue 678). Il ne porte que neuf lignes en haut, barrées ensemble de deux grands traits croisés ; le reste de la page est blanc et laisse voir le recto par transparence. En haut à droite, un trait oblique à crochet, peut-être un 1 ; pas de nombre de la série à l'encre. Ces lignes ne suivent pas la fin du recto (« $s$ une fonction de $\underline{r}$ seule. ») : elles commencent au milieu d'une énumération et semblent la fin d'un passage écrit ailleurs.}

\struck{$g_{1}, g_{2}, g_{3}$ \&c $g'_{1}, g'_{2}, g'_{3}$ \&c. $g^{2}_{1}, g^{2}_{2}, g^{2}_{3}$ \&c $\ldots\ldots$ $g^{\beta}_{1}, g^{\beta}_{2}, g^{\beta}_{3}$ \&c.
qui satisfont aux équations déterminées correspondantes aux
valeurs $1, 3, 5 \ldots\ldots 2\beta + 1$ du nombre $n$. $u_{1}, u_{2}, u_{3}$ \&c
$u'_{1}, u'_{2}, u'_{3}$ \&c. $u^{2}_{1}, u^{2}_{2}, u^{2}_{3}$ \&c $\ldots\ldots$ $u^{\beta}_{1}, u^{\beta}_{2}, u^{\beta}_{3}$ \&c, désignent
les valeurs de $\underline{u}$ qui correspondent a ces differentes
racines; $B^{0}$; $A, B$; $A', B'$; $A^{2}, B^{2}$; et $a, a_{2}, a_{3}$ \&c.
$a'_{1}, a'_{2}, a'_{3}$ \&c $\ldots\ldots$ $a^{\beta}_{1}, a^{\beta}_{2}, a^{\beta}_{3}$ \&c. sont des coefficiens
arbitraires.}
\note{Les indices supérieurs du deuxième groupe de chaque suite sont un petit trait, rendu par l'accent ($g'$, $u'$, $a'$, $A'$, $B'$), ceux des groupes suivants un 2 et un $\beta$. Le premier $a$ de la dernière suite et le couple « $A, B$ » n'ont pas d'indice visible.}

\page{680}

\note{En haut à droite, à l'encre, « 334 » ; au-dessus, séparé par un trait horizontal, un chiffre bouclé barré d'une croix, lu \uncertain{5} (il pourrait être un 8 ou un 7). Le feuillet est monté bas. Le texte suit la vue 678 (« $s$ une fonction de $\underline{r}$ seule ») ; le verso de la vue 679, barré, s'intercale entre les deux sans en faire partie.}

on aura $\frac{1}{r}\,\frac{du}{dr} = \text{sin}\,\beta\varphi\left(\beta r^{\beta-2}s + r^{\beta-1}\,\frac{ds}{dr}\right)$
\[ \frac{d^{2}u}{dr^{2}} = \text{sin.}\,\beta\varphi\left(\beta(\beta-1).\,r^{\beta-2}s + 2\beta r^{\beta-1}\,\frac{ds}{dr} + r^{\beta}\,\frac{d^{2}s}{dr^{2}}\right) \]
\[ \frac{1}{r^{2}}\,\frac{d^{2}u}{d\varphi^{2}} = -\beta^{2}\,\text{sin}\,\beta\varphi.\,r^{\beta-2}s \]
et la condition précédente deviendra:
\[ \frac{m}{k}\,s + \frac{d^{2}s}{dr^{2}} + \frac{(2\beta+1)}{r}\,\frac{ds}{dr} = 0 \]
\note{Dans la première ligne, le $r$ du second terme est noirci et l'exposant $\beta - 1$ est écrit presque sur la ligne. La nouvelle condition omet le facteur commun $\text{sin}\,\beta\varphi.\,r^{\beta}$, par lequel elle a été divisée.}

On choisira pour $\underline{s}$ une fonction de $\underline{r}$ qui satisfasse à cette
équation.
$g$ désignant la constante $\frac{m}{k}$ et $\underline{n}$ étant déterminée par la
condition $2\beta + 1 = n$, la fonction cherchée peut être exprimée
par la série suivante:
\[ s = 1 - \frac{gr^{2}}{2(n+1)} + \frac{g^{2}r^{4}}{2.4.(n+1)(n+3)} - \frac{g^{3}r^{6}}{2.4.6\,(n+1)(n+3)(n+5)} + \&c. \]
car on a: \quad $\frac{2\beta+1}{r}.\,\frac{ds}{dr} = -\,\frac{ng}{n+1} + \frac{ng^{2}r^{2}}{2(n+1)(n+3)} - \frac{ng^{3}r^{4}}{2.4\,(n+1)(n+3)(n+5)} + \&c.$
\[ \frac{d^{2}s}{dr^{2}} = -\,\frac{g}{n+1} + \frac{3g^{2}r^{2}}{2(n+1)(n+3)} - \frac{5g^{3}r^{4}}{2.4\,(n+1)(n+3)(n+5)} + \&c. \]
ainsi la somme des deux derniers termes de l'équation
sera égale au premier terme dela même équation pris
avec un signe contraire.
\note{Sous « cherchée », une petite tache d'encre. Le calcul est juste : la somme des deux lignes vaut $-g + \frac{g^{2}r^{2}}{2(n+1)} - \ldots$, c'est-à-dire $-gs$.}

Dans la valeur de $\underline{u}$ on peut \struck{également} mettre
\struck{$\text{sin.}\,\beta\varphi$} $\text{cos}\,\beta\varphi$ a la place de $\text{sin.}\,\beta\varphi$, parconséquent
$A$ et $B$ étant deux constants arbitraires
$e^{-gkt}\left(A\,\text{sin.}\,\beta\varphi + B\,\text{cos.}\,\beta\varphi\right)r^{\beta}s$ sera une valeur particulière
de $\underline{v}$.

L'état de la surface convexe du cylindre est
assujeti à une condition exprimée par l'équation
$hv + \frac{dv}{dr} = 0$, qui doit être satisfaite lorsque
\note{« constants » : ainsi au masculin sur le feuillet. Le $s$ final de « $r^{\beta}s$ » est récrit sur une tache. Dans la dernière équation, $h$ tient lieu du coefficient $\frac{h}{k}$, comme l'annonçait la vue 678. La page s'arrête au bas du feuillet sur « lorsque » ; la suite est hors de ce lot (vue 681 ou au-delà).}

\end{document}
